The red giant AFGL 3068 is located about $1.0\,\mathrm{kpc}$ from Earth. Its
images reveal a nearly Archimedean spiral with a constant arm spacing,
described by $r = a + b\phi$ where $a$ and $b$ are constants, and $\phi$ is
azimuthal angle in radian.

\ioaafig{2026_TH_T08-f1.pdf}

Such a spiral structure can be formed around a binary system. A red giant with
a fast stellar wind orbiting the binary center of mass produces a thin,
enhanced-density shell that expands radially at the wind velocity of $V_w$ in
the center of mass coordinate system. Its shapes on both the orbital plane
($Oxy$) and a perpendicular plane ($Oxz$) are shown below, where $\phi$ and
$\theta$ denote the azimuthal and polar angles, respectively. Let $r_p$ and
$V_p$ be the orbital radius and velocity of the red giant, respectively.

\ioaafig{2026_TH_T08-f2.pdf}

\medskip\noindent\textbf{Part 1 (10 points):}

\begin{description}
\item[(a)] (2 points) The on-sky arm spacing of AFGL 3068 is
  $2.29\,\mathrm{arcsec}$. Determine the constant $b$ in
  $\mathrm{au\,rad^{-1}}$.
\item[(b)] (4 points) At every point on the red giant trajectory, we consider
  only the wind moving radially outward with respect to the center of the
  binary system. Show that, on the orbital plane, the shell forms an
  Archimedean spiral structure with $a = r_p$ and $b = \dfrac{r_p V_w}{V_p}$.
\item[(c)] (2 points) The wind velocity of AFGL 3068 is
  $V_w = 14.0\,\mathrm{km\,s^{-1}}$. Calculate its orbital period $T$ in
  years.
\item[(d)] (2 points) The total mass of the binary system is
  $6.6\,\mathrm{M}_\odot$. Calculate the separation $d$ between the two stars
  in $\mathrm{au}$.
\end{description}

\medskip\noindent\textbf{Part 2 (15 points)}: Assume that the orbital plane is
inclined by an angle $i$ with respect to the sky plane. The $x$-axis is along
the intersection of the orbital and sky planes. $r_p$ is negligible compared
with the spiral size.

\begin{description}
\item[(a)] (6 points) Express the coordinates $x, y, z$ and velocities
  $(V_x, V_y, V_z)$ of a shell element in terms of the parameters
  $(\phi, \theta)$.
\item[(b)] (7 points) Show that all shell elements with line-of-sight velocity
  $V_{\mathrm{los}} = 0$ satisfy the equation of an ellipse on the sky plane
  with two semi-axes of $r\sin\theta$ and $r\sin\theta/\cos i$.
\item[(c)] (2 points) Recent CO line images at $V_{\mathrm{los}} = 0$ channel
  show that the spiral pattern of AFGL 3068 is slightly elongated relative to
  an Archimedean spiral. The elongated pattern is well described by ellipses
  with a ratio of $1.1$ between the two semi-axes. Determine the inclination
  angle $i$.
\end{description}
